% Source: https://artofproblemsolving.com/community/c3h1568400_mc_handbook_q
% Stable identity: aops-topic-1568400-q01
Let the three-digit number be $T$ and the two-digit number $t$, so $T=t+288$ with $10\le t\le98$, hence $298\le T\le386$.

\textbf{Hundreds digit.} If the hundreds digit of $T$ were $2$, then $T$ would be $298$ or $299$, which need a $9$. So $T=3XY$ with $X,Y$ from $\{0,2,7,8\}$, and
\[
t=T-288=300+\overline{XY}-288=\overline{XY}+12 .
\]
\textbf{The other digits.} The digits $X$, $Y$ and the two digits of $t$ must be $0$, $2$, $7$, $8$ in some order. Going through the twelve possible $\overline{XY}$ ($02,07,08,20,27,28,70,72,78,80,82,87$), the values of $\overline{XY}+12$ are $14,19,20,32,39,40,82,84,90,92,94,99$, and only $\overline{XY}=70$, $t=82$ uses exactly the remaining digits. So $T=370$, $t=82$, and the sum is
\[
370+82=\boxed{452}.
\]
(Check: $370-82=288$. The numbers do have to be found, but the casework is short once the hundreds digit is fixed.)
